% WHAT \showthe LEAVES BEHIND.
%
% \showthe prints what \the would have produced -- and that is ALL it
% consumes. The scan that reads a register number stops at the first
% token that cannot be part of it and puts that token back, and \showthe
% has no business with it: the reference names it under the answer,
%
%   > 0.
%   <to be read again>
%                      X
%
% and the X is read again afterwards, so `\hbox{\showthe\count0X}' holds
% an X. HSTeX took EVERYTHING the stack had gained back off it, which
% swallowed the terminator: the box came out empty, and in a formula the
% swallowed token was the `$' that should have closed it, so the rest of
% the run was one long recovery.
%
% A value \the gives is an inserted list and a token put back is not,
% which is what tells the two apart when the value is empty and no list
% was pushed for it at all.
%
% trip line 319 writes `\showbox22\kern3pt', and lines 273 and 317 ask
% \showthe\lastskip in the middle of a formula.
\catcode`\{=1 \catcode`\}=2 \catcode`\$=3
\font\rip=trip \rip
\tracingonline=1 \showboxbreadth=5 \showboxdepth=5
\message{[A the terminator is a letter, and is set]}
\setbox1=\hbox{\showthe\count0A}\showbox1
\message{[B the terminator ends a glue register's number]}
\setbox1=\hbox{\showthe\skip0A}\showbox1
\message{[C the terminator is the $ that closes a formula]}
\setbox0=\hbox{$A\showthe\count0$}
\message{[D a space after the number is eaten, and nothing is put back]}
\setbox1=\hbox{\showthe\count0 A}\showbox1
\message{[E nothing is put back after a quantity that reads no number]}
\setbox1=\hbox{\showthe\parindent A}\showbox1
\message{[done]}
\end
